\documentclass[a4paper,12pt]{article}

\usepackage[utf8]{inputenc}
\usepackage[T1]{fontenc}
\usepackage{lmodern}
\usepackage{graphicx}
\usepackage{hyperref}
\usepackage{listings}
\usepackage{xcolor}
\usepackage{enumitem}
\usepackage{tcolorbox}
\usepackage{float}

\hypersetup{
    colorlinks=true,
    linkcolor=blue,
    filecolor=magenta,
    urlcolor=cyan,
    pdftitle={Interaktywne uczenie asocjacyjne magistrali CAN},
    pdfauthor={CAN Simulator GUI},
    pdfsubject={Dokumentacja funkcji uczenia asocjacyjnego online},
    pdfkeywords={CAN, SocketCAN, nauczanie maszynowe, etykietowanie, interaktywne}
}

\lstset{
    basicstyle=\ttfamily\small,
    backgroundcolor=\color{gray!10},
    frame=single,
    breaklines=true,
    showstringspaces=false,
    language=Python,
    morekeywords={self,import,from,class,def,if,else,elif,for,while,return,True,False,None},
}

\title{Interaktywne uczenie asocjacyjne magistrali CAN\\Dokumentacja techniczna}
\author{Zespół CAN Simulator GUI}
\date{\today}

\begin{document}

\maketitle
\tableofcontents
\newpage

\section{Wprowadzenie}
Interaktywne uczenie asocjacyjne (ang.\ \emph{Interactive Associative Learning}) to nowa funkcja aplikacji CAN Simulator GUI umożliwiająca szybką identyfikację ramek CAN odpowiedzialnych za konkretne zjawisko fizyczne (np.\ włączenie świateł stopu, naciśnięcie pedału hamulca, otwarcie drzwi).

W przeciwieństwie do tradycyjnych metod (wyszukiwanie binarne, analiza logów \emph{post mortem}), użytkownik wskazuje momenty występowania zdarzenia w czasie rzeczywistym, a aplikacja automatycznie – na podstawie wielokrotnych obserwacji – wskazuje najbardziej prawdopodobną ramkę lub sekwencję ramek.

\section{Zasada działania}
\begin{enumerate}[noitemsep]
    \item \textbf{Tryb uczenia} – Użytkownik włącza tryb ,,Asocjacja CAN'' (zakładka \texttt{Uczenie asocjacyjne}).
    \item \textbf{Buforyzacja} – Aplikacja nasłuchuje ramek i przechowuje w buforze kołowym ostatnich $N$ sekund (domyślnie 10~s). Każda ramka zachowuje znacznik czasu, identyfikator, dane i flagę rozszerzonego ID.
    \item \textbf{Etykietowanie}:
        \begin{itemize}
            \item Gdy użytkownik \textbf{wciśnie} checkbox ,,Zdarzenie'' (lub klawisz \texttt{Ctrl+H}), rejestrowany jest czas $t_{\text{start}}$.
            \item Gdy użytkownik \textbf{puści} checkbox, rejestrowany jest czas $t_{\text{stop}}$.
            \item Ramki w przedziale $[t_{\text{start}}-\Delta t,\; t_{\text{stop}}+\Delta t]$ oznaczane są jako \textbf{pozytywne} (związane ze zdarzeniem). Pozostałe ramki w buforze to próbki negatywne.
        \end{itemize}
    \item \textbf{Analiza}:
        \begin{itemize}
            \item \emph{Heurystyka bajtowa} – dla każdego identyfikatora CAN i każdego bajtu sprawdzane jest, czy wartość bajtu jest stała w próbkach negatywnych, a zmienia się podczas zdarzenia. Jeśli tak, bajt i jego wartość stają się kandydatem.
            \item \emph{Klasyfikator online} (opcjonalny) – model \texttt{PassiveAggressiveClassifier} ze \texttt{scikit-learn} trenowany przyrostowo (\texttt{partial\_fit}) na każdej nowej próbce. Cechy wejściowe: ID, flaga extended, DLC oraz wartości 8 bajtów danych.
        \end{itemize}
    \item \textbf{Wynik} – lista kandydatów szeregowana malejąco według \emph{pewności} (średnia z heurystyki i modelu). Gdy pewność przekroczy próg (domyślnie 80\%), ramka jest uznawana za zidentyfikowaną.
    \item \textbf{Podświetlanie w Snifferze} – zidentyfikowane ramki są zaznaczane na żółto w zakładce ,,Sniffer CAN'', co ułatwia dalszą obserwację.
\end{enumerate}

\section{Interfejs użytkownika}

\subsection{Zakładka \texttt{Uczenie asocjacyjne}}

\begin{table}[H]
\centering
\begin{tabular}{|l|p{10cm}|}
\hline
\textbf{Element} & \textbf{Opis} \\
\hline
Checkbox ,,Zdarzenie'' & Przełącza stan zdarzenia (wciśnięty – rozpoczęcie, puszczony – zakończenie). Dostępny również globalnie przez \texttt{Ctrl+H}. \\
\hline
Przycisk ,,Start uczenia'' & Rozpoczyna buforowanie ramek i nasłuchiwanie. Wymaga aktywnego połączenia CAN. \\
\hline
Przycisk ,,Stop uczenia'' & Zatrzymuje buforowanie i analizę. \\
\hline
Przycisk ,,Wyczyść dane'' & Resetuje wszystkie zebrane dane, iteracje i tabelę wyników. \\
\hline
Tolerancja czasowa & Pole wyboru $\pm$X~ms – margines dodawany do okna zdarzenia (domyślnie $\pm200$~ms). \\
\hline
Tabela wyników & Kolumny: \texttt{CAN ID}, \texttt{Bajt}, \texttt{Wartość (zdarzenie)}, \texttt{Wartość (tło)}, \texttt{Pewność (\%)}. \\
\hline
Pasek postępu & Wizualizacja liczby zebranych iteracji (cel: 10). \\
\hline
Podgląd bufora & Wyświetla ostatnie 50 ramek z bufora w czasie rzeczywistym. \\
\hline
\end{tabular}
\caption{Elementy interfejsu zakładki \texttt{Uczenie asocjacyjne}.}
\end{table}

\subsection{Skróty klawiszowe}
\begin{itemize}[noitemsep]
    \item \texttt{Ctrl+H} – przełączenie stanu zdarzenia (działa globalnie, również gdy zakładka nie jest aktywna).
\end{itemize}

\section{Architektura}

\subsection{Struktura plików}
\begin{verbatim}
controllers/associative_controller.py   # Główna logika (bufor, etykietowanie, analiza, model)
gui/tabs/associative_tab.py             # Interfejs użytkownika (Tkinter)
tests/test_associative.py               # Testy integracyjne z vcan0
docs/associative_learning.tex           # Niniejszy dokument
\end{verbatim}

\subsection{Komponenty w \texttt{AssociativeController}}
\begin{itemize}[noitemsep]
    \item \texttt{CircularBuffer} – lista ramek z automatycznym usuwaniem starszych niż \texttt{max\_age} sekund.
    \item \texttt{toggle\_event()} – rozpoczyna/kończy zdarzenie; po zakończeniu wywołuje \texttt{\_label\_frames()} i \texttt{\_analyze()}.
    \item \texttt{\_extract\_features(rec)} – zamienia rekord ramki na wektor liczbowy dla klasyfikatora.
    \item \texttt{\_update\_classifier(pos, neg)} – uczy model \texttt{PassiveAggressiveClassifier} na bieżących danych.
    \item \texttt{get\_highlight\_ids(threshold)} – zwraca listę ID do podświetlenia w Snifferze.
\end{itemize}

\section{Testy}

\subsection{Test integracyjny}
Plik \texttt{tests/test\_associative.py} zawiera przypadki testowe:
\begin{enumerate}[noitemsep]
    \item \texttt{test\_event\_toggle\_counts\_iteration} – weryfikuje, czy przełączanie zdarzenia zwiększa licznik iteracji.
    \item \texttt{test\_buffer\_fills\_with\_messages} – wysyła ramki na \texttt{vcan0} i sprawdza, czy trafiają do bufora.
    \item \texttt{test\_labeling\_after\_event} – sprawdza oznaczenie ramek po zakończeniu zdarzenia.
    \item \texttt{test\_candidates\_generated} – symuluje zdarzenie z dedykowaną ramką i weryfikuje, że jest ona wysoko w rankingu kandydatów.
\end{enumerate}

\subsection{Uruchomienie testów}
\begin{lstlisting}[language=bash]
# Aktywuj wirtualne środowisko
source .venv/bin/activate

# Uruchom testy
python3 -m pytest tests/test_associative.py -v
\end{lstlisting}

\section{Konfiguracja}

\begin{table}[H]
\centering
\begin{tabular}{|l|l|l|}
\hline
\textbf{Parametr} & \textbf{Domyślnie} & \textbf{Opis} \\
\hline
\texttt{max\_age} & 10.0 & Maksymalny wiek ramki w buforze (sekundy) \\
\hline
\texttt{tolerance\_ms} & 200 & Margines tolerancji okna zdarzenia (ms) \\
\hline
\texttt{threshold} & 80.0 & Próg pewności do podświetlania (\%) \\
\hline
\end{tabular}
\caption{Parametry konfiguracyjne w \texttt{AssociativeController}.}
\end{table}

\section{Znane ograniczenia}

\begin{enumerate}[noitemsep]
    \item \textbf{Synchronizacja czasu} – opóźnienie między fizycznym zdarzeniem a reakcją użytkownika może wpływać na dokładność. Rozwiązaniem jest zwiększenie tolerancji czasowej.
    \item \textbf{Wielokrotne ID dla jednego zdarzenia} – system zwraca wszystkie korelujące ID, ale nie rozstrzyga, które są bezpośrednią przyczyną.
    \item \textbf{Szum danych} – ramki cykliczne o zmiennej wartości (RPM, prędkość) mogą być błędnie wskazywane jako kandydaci przy małej liczbie iteracji. Więcej iteracji redukuje ten efekt.
    \item \textbf{Klasyfikator online} – wymaga \texttt{scikit-learn}. Jeśli biblioteka nie jest dostępna, system działa wyłącznie na heurystyce bajtowej.
\end{enumerate}

\section{Przyszłe rozszerzenia}
\begin{itemize}[noitemsep]
    \item Filtrowanie wstępne – automatyczne odrzucanie ID o stałej częstotliwości i niewielkiej zmienności.
    \item Analiza sekwencji – wykrywanie ciągu ramek poprzedzających zdarzenie.
    \item Eksport wzorca – zapis znalezionej ramki/reguły do formatu umożliwiającego późniejszą automatyczną detekcję.
    \item Tryb pasywny – możliwość oznaczania zdarzeń na podstawie zewnętrznego sygnału (np.\ zmiana wartości konkretnej ramki).
\end{itemize}


\section{Wykrywanie sekwencji asocjacyjnych (Faza 9)}
Po zidentyfikowaniu głównej ramki (np. ID świateł stopu), użytkownik może kliknąć przycisk \texttt{Szukaj sekwencji}. Algorytm analizuje sąsiedztwo czasowe każdego wystąpienia ramki w buforze:

\begin{itemize}[noitemsep]
    \item Przeszukuje okno $\pm 0.5$ s wokół każdego wystąpienia głównego ID.
    \item Zlicza wystąpienia sąsiednich ID i ich kierunek (przed / po).
    \item Filtruje te, które występują w przynajmniej 80\% okien.
    \item Sortuje je według średnich odstępów czasu, tworząc uporządkowaną listę ID.
\end{itemize}

Wynik wyświetlany jest w kolumnie \texttt{Sekwencja} tabeli wyników. Wszystkie ID należące do sekwencji są podświetlane w Snifferze.

\section{Podsumowanie}
Funkcja interaktywnego uczenia asocjacyjnego stanowi praktyczne narzędzie diagnostyczne, pozwalające w kilka minut zidentyfikować nieznane ramki CAN odpowiedzialne za obserwowane zjawisko fizyczne. Dzięki integracji ze Snifferem i opcjonalnemu klasyfikatorowi online, proces identyfikacji jest szybki i intuicyjny.

\end{document}
